%==================================================================
%  COVER LETTER PREAMBLE  --  layout only, no content.
%  Loaded by src/lib/latex/compile.ts for documents of kind
%  "cover_letter". Keep in sync with preamble.tex where they share
%  a convention (no hyphenation, ragged right, empty page style).
%==================================================================
% The base size stays 10pt on purpose. With 11pt the class loads size11.clo,
% which asks for Computer Modern metrics at 10.95pt; the production image ships
% Latin Modern only, so that lookup fails before any text is set. The body size
% is raised after lmodern is loaded instead, which is why the line below exists.
\documentclass[a4paper,10pt]{article}

\usepackage[margin=2.2cm,top=2.0cm,bottom=2.0cm]{geometry}
\usepackage[T1]{fontenc}
\usepackage[utf8]{inputenc}
\usepackage{lmodern}

% Body text reads better one step above the CV. Latin Modern is scalable and
% already selected, so this costs nothing at compile time.
\makeatletter
\renewcommand{\normalsize}{\@setfontsize\normalsize{11}{13.6}}
\makeatother

\usepackage{microtype}
\usepackage{ragged2e}
\usepackage[hidelinks]{hyperref}

% Machine-readable text layer: ATS parsers and copy-paste.
\ifdefined\pdfgentounicode
  \input{glyphtounicode}
  \pdfgentounicode=1
\fi

\urlstyle{same}
\pagestyle{empty}
\raggedbottom
\setlength{\parindent}{0pt}
\setlength{\parskip}{7pt}
\setlength{\emergencystretch}{3em}
\RaggedRight
% No hyphenation, same as the CV: "colla-bora-tion" extracts as garbage.
\hyphenpenalty=10000
\exhyphenpenalty=10000

%---------------------------- macros ------------------------------
% Sender block: name, contact line, optional links line, then a hairline rule.
\newcommand{\letterhead}[3]{%
  {\LARGE\bfseries #1}\par\vspace{2pt}%
  {\small #2}\par
  \if\relax\detokenize{#3}\relax\else{\small #3}\par\fi
  \vspace{7pt}\noindent\rule{\linewidth}{0.5pt}\par
}

% Date, right aligned under the rule.
\newcommand{\letterdate}[1]{\vspace{2pt}\hfill{\small #1}\par\vspace{2pt}}

% Recipient block. The second argument is optional (role or team line).
\newcommand{\letterto}[2]{%
  \textbf{#1}\par
  \if\relax\detokenize{#2}\relax\else{\small #2}\par\fi
}

% Subject line: bold, no "Re:" theatre.
\newcommand{\lettersubject}[1]{\vspace{2pt}\textbf{#1}\par\vspace{2pt}}

% Salutation, kept on the same baseline flow as the body.
\newcommand{\letteropening}[1]{#1\par}

% Sign-off: the closing word, a signature gap, then the printed name.
\newcommand{\letterclosing}[2]{%
  \vspace{4pt}#1\par\vspace{16pt}\textbf{#2}\par
}
